%% kapa: exam=1 round=35 year=2024 period=2 form=A subject=accounting subject_ko=회계학 seq=20 number=60
%% source: Q-Net 감정평가사 기출문제 (cst00309) artlSeq=5213229 file=raw/1008/5213229/2교시 문제지 원본.pdf
%% conversion: kordoc 4.24.51 -> tools/gen1.py (auto); review=proofread-sonnet; flags=none
\begin{question}{60}
특수관계자 공시에 관한 설명으로 옳지 않은 것은?
\begin{choices}
  \item 보고기업에 유의적인 영향력이 있는 개인이나 그 개인의 가까운 가족은 보고기업의 특수관계자로 보며, 이 때 개인의 가까운 가족의 범위는 자녀 및 배우자로 한정한다.
  \item 지배기업과 종속기업 사이의 관계는 거래의 유무에 관계없이 공시한다.
  \item 특수관계자거래가 있는 경우, 재무제표에 미치는 특수관계의 잠재적 영향을 파악하는 데 필요한 거래, 채권ㆍ채무 잔액에 대한 정보뿐만 아니라 특수관계의 성격도 공시한다.
  \item 기업의 재무제표에 미치는 특수관계자거래의 영향을 파악하기 위하여 분리하여 공시할 필요가 있는 경우를 제외하고는 성격이 유사한 항목은 통합하여 공시할 수 있다.
  \item 지배기업과 최상위 지배자가 일반이용자가 이용할 수 있는 연결재무제표를 작성하지 않는 경우에는 일반이용자가 이용할 수 있는 연결재무제표를 작성하는 가장 가까운 상위의 지배기업의 명칭도 공시한다.
\end{choices}
\end{question}
\answer{1}
